%_____________________________________________________________________________________________
% Chapter 7: Conclusion and Future Work
%_____________________________________________________________________________________________

\chapter{Conclusion and Future Work}
\label{ch:conclusion}

\section{Conclusion}
This project set out to build an ML model drift early-warning system that warns about degradation, not merely drift, before delayed labels can confirm it. We built a benchmark that re-samples real tabular data so that whether a shift harms the model is measured from true labels, and we evaluated ten versions of the system, five of them with pre-registered hypotheses, frozen thresholds and one-time runs on unseen datasets.

The results are mixed, and we report them as they are:
\begin{itemize}
\item The original design, a learned risk calibrator over raw drift scores, failed all its hypotheses, because thresholds in dataset-specific units did not transfer across datasets.
\item A stream-normalised \textbf{drift watch} gave early warning of harmful drift within its false-alert budget on benchmark streams: on six fresh datasets it warned before the accuracy drop on 87\% of harmful gradual shifts, a median of six batches ahead. It cannot tell harmful from harmless drift, and it failed on natural streams where drift had begun before monitoring.
\item An importance-weighted \textbf{harm tier} reduced alarms on harmless shifts from 97\% to 34\% and separated harmful from harmless shifts with an AUROC of 0.78, short of the 0.80 target but above confidence-based baselines on the same data. On a later set of datasets its false-alert rate exceeded the budget.
\item Nonlinear estimators, ensembles and stricter alerting rules did not improve on it.
\end{itemize}
Without labels, our best estimator separated harmful from harmless shifts only partly, and it needed a labeled reference set. The practical lessons are to express alert scores in units that are comparable across streams, to validate thresholds on data from other sources, to separate a sensitive drift tier from a specific harm tier, and to keep a labeled reference sample.

Beyond the research findings, the project delivers a working and reproducible product: the monitoring library, frozen configurations and verdict scripts, a dashboard that replays monitored streams with delayed labels, 85 unit tests with continuous integration, and a research paper.

\section{Future Work}
\begin{enumerate}
\item \textbf{Better burn-in handling.} Detect whether the burn-in itself is stable, or use a reference period chosen with labels, so that stream normalisation also works when drift has begun before monitoring.
\item \textbf{Calibrated harm alerting.} Control the false-alert rate over a whole monitoring period (for example with sequential testing) instead of per batch, so that per-episode alarm targets become achievable.
\item \textbf{Richer density-ratio estimation.} Combine importance weighting with uncertainty estimates, and test other models than logistic regression as the monitored model.
\item \textbf{Label changes.} Use the few labels that arrive early (partial labeling) to detect changes in $P(Y\mid X)$, which no input-only signal can see.
\item \textbf{More data types.} Extend the benchmark to multi-class, text and image data, and to real temporal streams with documented drift.
\item \textbf{Deployment.} Package the two tiers as a service that reads batches from a production pipeline, with the dashboard as its operator view.
\item \textbf{Public pre-registration.} Register future experiments on a public registry before running them.
\end{enumerate}
